\documentclass[10pt,a4paper]{amsart}
\usepackage[margin=19mm]{geometry}
\usepackage{amsmath,amssymb,mathtools}
\usepackage[scr=rsfs,cal=euler]{mathalfa}
\usepackage{xfrac}
\usepackage[unicode,hidelinks,bookmarks=false]{hyperref}
\newcommand{\ie}{i.e.\ }
\newcommand{\cf}{cf.\ }
\newcommand{\resp}{respectively\ }
\newcommand{\Subsection}{\subsection}
\providecommand{\nobreakdash}{\nobreak\hbox{-}}
\setlength{\emergencystretch}{2em}
\multlinegap0pt
\usepackage{bm}
\usepackage{bbm}
\usepackage{dsfont}
\usepackage{xspace}
\usepackage{cancel}
\usepackage{mathtools}
\usepackage{cleveref}
\usepackage{amsthm}
\makeatletter
\@ifundefined{theorem}{\newtheorem{theorem}{Theorem}}{}
\@ifundefined{lemma}{\newtheorem{lemma}[theorem]{Lemma}}{}
\@ifundefined{proposition}{\newtheorem{proposition}[theorem]{Proposition}}{}
\makeatother
\def\GG{\mathbb{G}}
\def\Z{\mathbb Z}
\def\mc{\mathcal}
\def\tsc{{\rm{\bf{fsc}}}}
\def\t{\tau}
\title[Profinite rigidity for free-by-cyclic groups with centre]{Profinite rigidity for free-by-cyclic groups with centre}
\author{Martin R. Bridson, Pawe{\l} Piwek}
\date{Source: arXiv:2409.20513v1 (CC BY 4.0)}
\begin{document}
\maketitle

\noindent\emph{Source and licence.} Profinite rigidity for free-by-cyclic groups with centre. Version arXiv:2409.20513v1 (CC BY 4.0), distributed under the Creative Commons Attribution 4.0 International licence, as recorded in the arXiv licence metadata for this exact version and shown on the abstract page https://arxiv.org/abs/2409.20513v1. Full rights notice: licenses/arxiv-2409.20513-CC-BY-4.0.txt.\par\smallskip

\noindent\textit{Editorial scope.} Counted unit: the Proposition whose source label is \texttt{associated\_poset\_is\_torsion\_subgroups} in the article (source0.tex lines 1893--1903, source label \texttt{associated\_poset\_is\_torsion\_subgroups}), delivered together with the article's own complete original proof of it (source0.tex lines 1905--1945, one \texttt{proof} environment of 41 lines). No mathematics is written, repaired or re-derived: statement and proof are the article's own text line for line. Required context retained uncounted: which\_graphs (lines 1139--1157); l:sliding-the-same (lines 1845--1851). Every definition, display and label of those lines is retained unchanged. Omitted from this package and not redistributed: 1--1138, 1158--1844, 1852--1892, 1904--1904, 1946--3023. Nothing inside a retained excerpt is omitted or altered: every retained body is delivered exactly as the article's lines are written, so no omission receipt of any kind accompanies this package. Citations: the retained excerpts carry no citation command at all, so no bibliography entry of the article is transcribed and no reference is redistributed. Reference adaptation: none was needed and none was made; every reference inside the delivered unit resolves inside it, so no unresolved reference or citation is produced. Printed slips kept literal: no printed word, symbol, hypothesis or number of the counted statement or of its proof was altered. Layout and class: the article's own class is \texttt{amsart} with packages including amssymb, amsfonts, amsmath, mathabx, color, tikz, subfigure, fontenc, inputenc, amsthm, thmtools, thm-restate, centernot, babel; the wrapper is the lane's pinned \texttt{amsart} editorial wrapper at 10pt on A4, which restates the theorem-like environments the retained text needs and adds the article's own macro definitions that the retained text needs (\texttt{\textbackslash GG}, \texttt{\textbackslash Z}, \texttt{\textbackslash mc}, \texttt{\textbackslash t}, \texttt{\textbackslash tsc}), with extra wrapper packages (bm, bbm, dsfont, xspace, cancel, mathtools, cleveref). The delivered numbering is the unit's own. Assets: no retained span contains a figure environment, a graphics-inclusion command, a PDF-page inclusion, a table or a TikZ picture; no image file is copied into this package, the package declares no asset dependency and no image file is redistributed with it. Licence: version arXiv:2409.20513v1 of this article is distributed under the Creative Commons Attribution 4.0 International licence, as recorded in the arXiv licence metadata for this exact version and shown on the abstract page https://arxiv.org/abs/2409.20513v1. A full rights notice is delivered at licenses/arxiv-2409.20513-CC-BY-4.0.txt. Characters: every retained excerpt is pure ASCII, so the delivered engine, XeLaTeX with its default Latin Modern text font, reports no missing character.
\begin{proposition}
\label{which_graphs} 
	A finitely generated group $G$ is free-by-(finite cyclic)
	if and only if there exists 
	a finite graph of finite cyclic groups $\GG(V,E)$ with 
	$G_v\cong \Z/d_v\ (e\in V)$ and 
	$G_e\cong \Z/d_e\ (e\in E)$ so that 
	\begin{enumerate}
		\item $G\cong\pi_1\GG$,
		\item $d_e$ divides $d_{\tau(e)}$ for all $e\in E$,
		\item there are generators $t_v$ for $G_v\ (v\in V)$ 
		and $t_e$ for $G_e=G_{\-e}\ (e\in E)$ so that 
		$$\iota_e(t_e) = t_{\t(e)}^{\delta_e}$$ 
		for all $e\in E$, where $\delta_e = d_{\tau(e)}/d_e$. 
	\end{enumerate} 
	Moreover, one can choose generators with these properties 
	for the local groups in any decomposition of $G$ 
	as a graph of finite groups. 
\end{proposition}

\begin{lemma}\label{l:sliding-the-same}
	There is a 1-1 correspondence  $\GG \leftrightarrow \mc{G}$
	between the finite graphs of finite groups described in Proposition \ref{which_graphs} and
	the set of finite numbered graphs. If $\GG_i$ corresponds to $\mc{G}_i$ for $i=1,2$, then
	there is a slide move transforming $\GG_1$ into $\GG_2$ if and only if there is a slide move
	transforming $\mc{G}_1$ into $\mc{G}_2$.
\end{lemma}

\begin{proposition}
\label{associated_poset_is_torsion_subgroups}
	Let $G$ be a finitely generated free-by-(finite cyclic) group, 
	let $\GG$ be a reduced graph of finite groups 
	with $G\cong \pi_1\GG$, 
	and let $\mc{G}$ be the numbered graph that encodes $\GG$ 
	(as in \Cref{l:sliding-the-same}).  
	
	Then, there is a data-preserving isomorphism 
	$\Psi: \tsc(G)\to \mathsf{P}_\mathcal{G}$. 
\end{proposition}

\begin{proof}
	Consider the action of $G=\pi_1\GG$ on
	the universal covering $T$ of $\GG$, which is the Bass-Serre tree of the splitting.
	The fixed-point set ${\rm{Fix}}(S)\subseteq T$ of each finite subgroup $S<G$
	is connected and non-empty, and its image in $\mc{G}=G\backslash T$, which we call  $X_S$, is the same as the image of
	${\rm{Fix}}(S^g)$ for every conjugate $S^g=g^{-1}Sg$. All of the edges and vertices in $X_S$
	have a copy of $S$ in the local group, and hence the numbers labelling the edges and vertices of $X_S\subseteq\mc{G}$
	are all divisible by $k=|S|$; in other words $X_S\subseteq\mc{G}_k$. In fact,
	$X_S$ is a connected component of $\mc{G}_k$: for each vertex $v\in X_S$ and each
	edge $e$ with $\tau(e)=v$, if $e$ is in $\mc{G}_k$ then $k$ divides
	the label on $e$, which means that  $G_e<G_v$ contains a subgroup of cardinality $k$, and the only such subgroup in  
	$G_v$ (which is cyclic) is $S$, hence $S \leq G_e$ and $e\subset X_S$.  Thus $X_S=X_{S^g}$  is one of the elements
	of level $|S|$ in  $\mathsf{P}_\mathcal{G}$ and we can define $\Psi([S]) := X_S$.
	Note that $\Psi$ is injective, because each of the local groups  in $\GG$ contains a unique subgroup
	of any given cardinality, so $X_{S_1}$ and $X_{S_2}$ are disjoint if $S_1$ and $S_2$ are non-conjugate subgroups
	of the same size.
	
	To see that $\Psi$ is surjective, consider an element $a_{i,k}\in \mathsf{P}_\mathcal{G}$. By definition,
	$a_{i,k}$ is a subgraph of $\mc{G}_k$, so the local group at each vertex $v$ of $a_{i,k}$ contains a subgroup of
	cardinality $k$. If this subgroup is $S$, then $v\in X_S \cap a_{i,k}$. 
	But $a_{i,k}$ and $X_S$ are connected components of $\mc{G}_k$, so having non-trivial intersection forces $a_{i,k}=X_S$,
	in other words $a_{i,k} = \Psi([S])$. Thus $\Psi$ is a bijection.
	
	If $S_1 \leq S_2$, then ${\rm{Fix}}(S_1)\supseteq {\rm{Fix}}(S_2)$ 
	and hence $X_{S_1}\supseteq X_{S_2}$.
	Moreover, $|S_1|$ divides $|S_2|$. 
	Thus $S_1 \leq S_2$ implies $\Psi([S_1]) \preceq \Psi([S_2])$,
	so $\Psi$ is an isomorphism of posets.
	
	It remains to prove that $\Psi$ preserves tags, i.e. that the first Betti number of $C_G(S)$ is equal to the
	first Betti number of $X_S = \Psi([S])$. For this, the first thing to observe is that 
	while $N_G(S)$ preserves the fixed-point set ${\rm{Fix}}(S)$ in the Bass-Serre tree $T$,
	the elements of $G$ that are not in $N_G(S)$ move ${\rm{Fix}}(S)$  off itself entirely. 
	Indeed,  if $v$ and $gv$ are both vertices of ${\rm{Fix}}(S)$, then $S$ and $S^g$ are both subgroups of the
	stabiliser of $v$, which being cyclic has only one subgroup of size $|S|$, so $S=S^g$ and $g\in N_G(S)$.
	It follows that $N_G(S)\backslash{\rm{Fix}}(S)$ injects into $G\backslash T = \mathcal{G}$ with image $X_S$.
	Thus $N_G(S)$ is the fundamental group of a graph of finite groups with underlying graph $X_S$,
	and therefore $N_G(S)$ has  the same first Betti number as $X_S$. 
	The final thing to observe is  $N_G(S) = C_G(S)$, because in any (torsion-free)-by-abelian group,
	the normaliser of a finite subgroup  is the centraliser of that subgroup. 
\end{proof}

\end{document}
